% Table 3: Curriculum Structure
\begin{table}[htbp]
\centering
\caption{Six-unit curriculum structure providing research infrastructure and progressive learning path. Units 1-4 are domain-specific (transferable), Units 5-6 are domain-agnostic (reusable across scientific disciplines).}
\label{tab:curriculum}
\scriptsize
\begin{tabular}{@{}clcp{4cm}@{}}
\toprule
\textbf{Unit} & \textbf{Topic} & \textbf{Notebooks} & \textbf{Key Technologies} \\
\midrule
1 & Nanoscale Modeling & 1 & ASE, RDKit, NGLView \\
2 & Molecular Simulation (MD, DFT) & 2 & ASE/EMT, OpenMM \\
3 & Machine Learning for Nanomaterials & 4 & scikit-learn, PyTorch \\
4 & Applied AI in Nanotechnology & 2 & Materials Project API, Computer Vision \\
\rowcolor{blue!10}
5 & Multi-Agent Systems & 10 & LangGraph, CrewAI, ADK, smolagents, Graph RAG, MCP \\
\rowcolor{blue!10}
6 & Integration Project & 6 & Full pipeline, FastAPI, Docker \\
\midrule
& \textbf{Total} & \textbf{25} & \\
\bottomrule
\end{tabular}

\vspace{0.15cm}
\scriptsize
\textit{Note:} Highlighted rows (Units 5-6) provide domain-agnostic infrastructure reusable for any computational science domain.
\end{table}
